\documentclass[a4paper, 11pt]{article}

\usepackage{graphicx, geometry, sectsty, marginnote, tgpagella}
\usepackage[semibold]{ebgaramond}
\usepackage[usenames, dvipsnames]{xcolor}
\usepackage[bookmarks, colorlinks, breaklinks]{hyperref}

\setlength{\parindent}{0pt}
\pagenumbering{gobble}

\geometry{
	paper=a4paper,
	top=3.5cm,
	bottom=2.5cm,
	left=4.5cm,
	right=3.5cm,
	headheight=0.75cm,
	footskip=1cm,
	headsep=0.75cm,
}

\sectionfont{\fontsize{15pt}{18pt}\selectfont}
\subsectionfont{\mdseries\scshape\normalsize}
\subsubsectionfont{\mdseries\upshape\bfseries\normalsize}

\newcommand{\years}[1]{\marginnote{\small #1}}
\renewcommand*{\raggedleftmarginnote}{}
\setlength{\marginparsep}{20pt}
\reversemarginpar

\hypersetup{
	linkcolor=blue,
	citecolor=blue,
	filecolor=black,
	urlcolor=MidnightBlue
}
\hyphenpenalty=10000

\begin{document}

\section*{\LARGE K Tejaswi Narayana}
3rd Year BS-MS \\
IISER Pune, \\
Dr. Homi Bhabha Road, Pune 411008. \\
Email: \href{mailto:k.tejaswinarayana@students.iiserpune.ac.in}{k.tejaswinarayana@students.iiserpune.ac.in}

\section*{Education}
\years{2024-2029 \\ (Expected)} {\bf Bachelor of Science, Master of Science} (BS-MS) in Mathematics\\
Indian Institute of Science Education and Research Pune. \\

%\years{2024} {\bf II PUC} (Class 12th) - {\bf 96.33\%} \\
%Sharada Pre-University College, Kodialbail, Mangaluru. \\

%\years{2022} {\bf SSLC} (Class 10th) - {\bf 97.92\%} \\
%Shri Dharmasthala Manjunatheshwara High School, Dharmasthala.

\section*{Project Experience}

\years{Summer 2026 (Ongoing)} {\bf Introduction to Lie Group Theory}, with Prof. Anupam Singh, IISER Pune. \\
\href{https://github.com/tejaswi-kodiadka/notes/blob/main/matrix-groups/matrix-groups.pdf}{Project Report}. (In Progress) \\\\

\years{Summer 2025} {\bf Basics of Affine and Projective Geometry}, with Prof. Steven Spallone, IISER Pune.

Studied some of the basic fundamental results in affine and projective geometry with the help of {\it Geometry} by M. Audin,
and {\it Geometry} by Brannan, Esplen, and Gray. \\
\href{https://github.com/tejaswi-kodiadka/notes/blob/main/geometry/geometry.pdf}{Project Report}. \\

\section*{Scholarships}
\years{2024-Present} {\bf INSPIRE-SHE Fellowship} \\
Awarded by the Department of Science and Technology, Govt. of India to the top 1\% of the Class XII performers.

\section*{Extracurricular}
\years{2025-2026} {\bf Chess Coordinator}, Sports Club \\
Organized regular sessions, events, and tournaments throughout the academic year. \\

\years{2026-2027} {\bf Question Making Coordinator for Mathematics}, Mimamsa \\
Mimamsa is a prestigious annual nation-wide science competition for undergraduate students, organized by
the students of IISER Pune. The problems in the competition are based on concepts usually taught in highschool,
but require critical and creative thinking to solve.

\end{document}
